\documentclass{article}

% ---------------------------------------------------------------------------
% CS2810 Embodied AI (Spring 2026)
% Combined project report: Goalkeeper + Shooter
% Template updated to the HW4 NeurIPS preprint style.
% Build: pdflatex main.tex  (run twice for references/links)
% ---------------------------------------------------------------------------

\usepackage[preprint,nonatbib]{neurips_2026}

\usepackage[utf8]{inputenc}
\usepackage[T1]{fontenc}
\usepackage[strings]{underscore} % allow bare _ in text and \texttt
\usepackage{url}
\usepackage{amsfonts}
\usepackage{nicefrac}
\usepackage{microtype}
\usepackage{amsmath}
\usepackage{booktabs}
\usepackage{array}
\usepackage{multirow}
\usepackage{graphicx}
\usepackage{xcolor}
\usepackage{caption}
\usepackage{enumitem}
\usepackage{listings}
\usepackage[colorlinks=true,linkcolor=blue,citecolor=blue,urlcolor=blue]{hyperref}

\captionsetup{font=small,labelfont=bf}
\setlist[itemize]{nosep,leftmargin=*}
\setlist[enumerate]{nosep,leftmargin=*}
\setlength{\tabcolsep}{4pt}
\renewcommand{\arraystretch}{0.95}

\lstset{
  basicstyle=\ttfamily\small,
  breaklines=true,
  columns=fullflexible,
  keepspaces=true,
  frame=single,
  framesep=4pt,
  xleftmargin=2pt,
  aboveskip=6pt,belowskip=6pt,
}

\newcommand{\code}[1]{\texttt{#1}}
\newcommand{\TODO}[1]{\textcolor{red}{\textbf{[TODO: #1]}}}

\title{\textbf{Learning Shooting and Goalkeeping Skills\\
for the Unitree G1 Humanoid}}
\author{Team 12\\ % member names to be added before submission
CS2810 Embodied AI, Spring 2026}
\date{}

\begin{document}
\maketitle

\begin{abstract}
\noindent
We develop reinforcement-learning soccer skills for the 29-DoF Unitree G1
humanoid in the mjlab MuJoCo simulator, covering the project's two tasks:
\emph{shooting} (approach and kick a stationary ball into the goal) and
\emph{goalkeeping} (intercept a ball launched on a parabolic trajectory). For
the shooter we build a six-stage progressive training pipeline around a shared
recurrent policy, reaching a \textbf{99.96\%} mean Phase-1 success rate
(Phase-1 score $\approx 29.9/30$) while maintaining high shot speed and
directional consistency. For tournament play, we deploy the shooter behind a
multi-strategy API server whose strongest mode selects motions and shot targets
online according to the goalkeeper's observed position. For the goalkeeper, we
organize the controller around an explicit two-phase design: a \emph{prepare}
phase that stabilizes the robot into a ready-to-dive stance before the shot is
fully committed, and an \emph{active} phase that executes the actual
interception once the incoming trajectory is clear. On top of this structure,
we use a staged pipeline that distils a diving prior, repairs difficult saves,
and specializes interception behavior by ball region. The resulting Phase-1
goalkeeper reaches a \textbf{94.67\%} block rate over the official
$150$-trial three-seed evaluation (Phase-1 score $22.0/30$). For both tasks,
we also describe how the competition server reconstructs policy observations
directly from raw MuJoCo state so that deployment matches training-time
observation construction.
\end{abstract}

\section{Shooter}\label{part:shooter}

\subsection{Phase 1 Results}\label{sec:sh-phase1}

\paragraph{Evaluation setup.} We use the official Phase-1 shooting evaluation
and supplement it with repeated large-scale evaluation over three fixed seeds to
reduce variance and confirm stability.

\begin{table}[h]
\centering\small
\caption{Shooter Phase-1 quantitative results.}
\label{tab:sh-phase1-summary}
\begin{tabular}{@{}lc@{}}
\toprule
Metric & Value \\
\midrule
Seed 42 success rate & $99.96\%$ \\
Seed 2810 success rate & $99.95\%$ \\
Seed 202606 success rate & $99.97\%$ \\
Mean success rate & $\mathbf{99.96\%}$ \\
Phase-1 score & $\mathbf{29.9/30}$ \\
Mean kick speed & $13.03$\,m/s \\
Kick-direction cosine similarity & $0.9787 \pm 0.0191$ \\
\bottomrule
\end{tabular}
\end{table}

\paragraph{Analysis.} The Phase-1 shooter is already close to saturation:
misses are extremely rare, and the small gap from full marks is due mainly to
the grading rule magnifying the tiny gap from $100\%$. The main remaining
weakness is far-corner precision under the hardest targets, not basic kick
generation.

\subsection{Phase 2 Design \& Approach}

\paragraph{Phase-2 approach.} For Phase 2, we refine the strong Phase-1
single-agent shooter for tournament play through competition-frame alignment,
motion selection, goalkeeper-aware target planning, and API-side deployment
logic. The objective is to improve robustness against a reacting goalkeeper
while preserving the high-accuracy, high-speed kicking behavior established in
Phase 1.

\subsubsection{Overview}\label{sec:sh-overview}

\textbf{Task.} The G1 humanoid (29 DoF) performs a penalty-spot shot in MuJoCo:
approach a stationary ball and kick it into the goal. This is a
perception-guided motor skill requiring a motion-tracking prior and
ball-trajectory control. Table~\ref{tab:sh-deliverables} summarizes the
deliverables.

\begin{table}[h]
\centering\small
\caption{Shooter deliverables.}
\label{tab:sh-deliverables}
\begin{tabular}{@{}lp{9.5cm}@{}}
\toprule
Artifact & Description \\
\midrule
Progressive training pipeline & A six-stage curriculum that starts from motion imitation and gradually shifts toward competition-style ball control and scoring. \\
Final shooter policy & A recurrent policy that preserves motion consistency while reacting to the ball and shot destination online. \\
Competition policy layer & A multi-strategy API server that supports random shooting, best-motion selection, and goalkeeper-aware target planning. \\
Evaluation analysis & Motion-level accuracy and speed analysis used to identify which kicks remain reliable under different target regions. \\
\bottomrule
\end{tabular}
\end{table}

\subsubsection{Training pipeline: six-stage progressive training}\label{sec:sh-pipeline}

Following the progressive perception--action framework of PAiD~\cite{kong}, we
design six training stages. All stages share the same recurrent policy
structure; what changes across stages is the task emphasis. The first stage
establishes a stable motion prior by imitating kick trajectories. The next
stage adds ball interaction, so the agent learns not just to reproduce motion
clips but to make controlled contact with the ball. The middle stages then
shift the objective toward scoring quality, explicitly favoring shots that are
both accurate and fast enough to remain threatening in a competitive setting.
Finally, the last stages align training fully with the competition world frame,
so that the learned behavior transfers directly to the tournament setup rather
than relying on an implicit train-to-deploy coordinate conversion.

\paragraph{Design conclusion.} The main lesson from this curriculum is that
high-quality kicking is easier to obtain by \emph{re-weighting the objective
across stages} than by searching for one perfect reward from the start. The
early motion prior stabilizes learning; the later competition-aligned stages
remove the remaining domain gap and produce the final near-saturated Phase-1
performance.

\subsubsection{Model architecture and observation space}\label{sec:sh-arch}

\paragraph{Architecture and observations.} The final shooter uses a recurrent
actor--critic architecture because successful kicking is fundamentally
sequential: approach, planting, swing timing, and shot redirection must remain
consistent over multiple control steps. On the actor side, the observation is
organized around three complementary sources of information: motion-reference
state to preserve mechanically plausible kicks, proprioceptive state to
stabilize the robot online, and task state describing the ball and desired shot
destination in the robot's local frame. During training, the critic is given
additional privileged information for value estimation, but deployment relies
only on the actor-side observation reconstructed from raw simulator state.

\paragraph{Design conclusion.} Three conclusions follow from this design.
First, recurrent state is essential: when the shot target is updated online, a
memoryless controller is much more likely to break the kick sequence. Second,
the observation should remain compact but structured, since mixing motion,
proprioception, and task context is enough to support strong performance without
needing a large hand-engineered state description. Third, exact deployment
consistency matters more than aggressive robustness tricks at the final stage:
once training is aligned to the competition frame, reducing observation and
environment randomization improves transfer because it removes the dominant
source of mismatch between training and tournament execution.

\subsubsection{Reward design evolution}\label{sec:sh-reward}

To make the curriculum progression explicit, Table~\ref{tab:sh-reward-evolution}
summarizes how the training objective shifts across stages.

\begin{table}[h]
\centering\footnotesize
\caption{Core reward-term evolution across stages.}
\label{tab:sh-reward-evolution}
\setlength{\tabcolsep}{4pt}
\begin{tabular}{@{}lccccccl@{}}
\toprule
Reward term & S1 & S2 & S3 & S4 & S5 & S6 & Function \\
\midrule
\code{track_body_pos}     & 1.0 & 1.0 & 0.5 & 0.5 & \textbf{0.05} & 0.05 & body pose tracking \\
\code{track_body_ori}     & 1.0 & 1.0 & 0.5 & 0.5 & \textbf{0.05} & 0.05 & body orientation tracking \\
\code{track_anchor_ori}   & 1.0 & 1.0 & 0.5 & 0.5 & \textbf{0.05} & 0.05 & anchor orientation tracking \\
\code{track_body_lin_vel} & 1.0 & 1.0 & 0.5 & 0.5 & \textbf{0.05} & 0.05 & body linear-velocity tracking \\
\code{track_body_ang_vel} & 1.0 & 1.0 & 0.5 & 0.5 & \textbf{0.05} & 0.05 & body angular-velocity tracking \\
\code{track_foot_pos}     & --- & 1.0 & 0.5 & 0.5 & \textbf{0.05} & 0.05 & foot position tracking \\
\textbf{tracking total}   & \textbf{6.0} & 5.0 & 3.0 & 3.0 & \textbf{0.3} & \textbf{0.3} & \\
\code{contact}            & --- & \textbf{50} & 50 & 50 & 50 & 50 & foot--ball contact force \\
\code{sideways_kick}      & --- & \textbf{50} & 50 & 50 & 50 & 50 & lateral kick toward goal \\
\code{ball_vel_align}     & --- & 30 & 30 & 30 & 30 & \textbf{20} & ball velocity aligned to goal \\
\code{ball_speed}         & --- & 10 & \textbf{15} & \textbf{20} & 20 & \textbf{40} & ball horizontal speed \\
\code{goal_accuracy}      & --- & --- & \textbf{10} & 10 & \textbf{30} & 20 & goal-plane crossing accuracy \\
\code{goal_miss}          & --- & --- & $-5$ & $-20$ & \textbf{$-50$} & $-50$ & miss penalty \\
\code{z_speed}            & --- & 0 & 0 & 0 & $-2$ & \textbf{$-5$} & vertical-speed penalty \\
\code{ball_oob}           & --- & --- & --- & --- & $-30$ & $-30$ & out-of-bounds penalty \\
\bottomrule
\end{tabular}
\end{table}

The reward design follows the same principle as the curriculum itself: early
training emphasizes \emph{stability and motion quality}, while later training
shifts toward \emph{shot outcome}. At the beginning, imitation-style tracking is
necessary because the humanoid first needs a coherent kicking posture and
timing. Once that prior is established, the reward increasingly favors
meaningful ball contact, goal-directed velocity, and eventual scoring outcome.
By the final stages, the optimization target is no longer ``look like a kick''
but ``produce a fast, accurate, competition-valid shot.''

\paragraph{Design conclusion.} The most important lesson is that successful
humanoid shooting requires a \emph{controlled handoff} from motion fidelity to
task utility. If tracking remains too dominant, the robot reproduces kicks that
look plausible but are not competitive. If task rewards dominate too early, the
robot loses motion quality and becomes unstable. The staged schedule solves this
by shifting the balance gradually rather than abruptly.

\subsubsection{Tournament strategy and API server}\label{sec:sh-phase2}

\paragraph{API server.} The shooter is deployed behind the standard
competition API. A key design choice is that the server reconstructs the full
policy observation from raw simulator state rather than relying on any shared
feature interface. This keeps the learned policy compatible with the tournament
protocol while preserving the exact observation structure used during training.

\paragraph{Observation computation.} The server combines proprioception, ball
state, motion-reference state, and the selected destination into the same
observation layout used by the recurrent policy during training. This matters in
practice because the shooter is sensitive not only to \emph{what} target is
chosen, but also to \emph{when} that target enters the recurrent control loop.

\paragraph{Three shooting strategies.}
\begin{table}[h]
\centering\small
\caption{Shooting strategies.}
\begin{tabular}{@{}llp{6.4cm}l@{}}
\toprule
Strategy & Trigger & Behavior & Latency \\
\midrule
\code{uniform} & \code{/reset} & random motion (incl. left leg) $+$ random target $y\in[-1.4,1.4]$, locked immediately & none \\
\code{best-motion} & \code{/reset} & random target $+$ \code{_select_motion_for_dest()} picks the best right-leg motion, locked immediately & none \\
\code{gk-aware} & per-frame \code{/act} & observe the GK for 5--16 frames, then lock the far-side target $+$ best motion & 5--16 frames \\
\bottomrule
\end{tabular}
\end{table}

\paragraph{The \code{gk-aware} strategy.}
This is the default and most capable strategy, in three phases.

\textbf{Observation phase} (\code{shoot_planned["done"]==False}): each \code{/act}
appends the GK root-$y$ to \code{gk_y_history}; \code{_should_lock()} applies a
dual condition -- \emph{behavioral trigger} (GK lateral offset $\geq0.3$\,m for a
run of \code{min_observe}$=5$ frames $\to$ early lock) and \emph{hard deadline}
(reaching \code{deadline_pct}$=60\%$ of the kick frame $\to$ force lock).
Crucially we do \textbf{not} reset the policy on lock: locking only changes the
\code{dest_body} observation term and leaves the LSTM hidden state intact, because
the policy was trained on a continuous motion stream and jumping to frame 16 with
an empty hidden state causes an out-of-distribution collapse.

\textbf{Target planning} (\code{_plan_destination}): if the GK is left
($y<-0.3$), shoot right \code{random(0.3,1.4)}; if right ($y>+0.3$), shoot left
\code{random(-1.4,-0.3)}; if central ($|y|\leq0.3$), shoot the full goal width
\code{random(-1.4,1.4)}. The random component within the chosen half adds
unpredictability.

\textbf{Execution:} \code{_select_motion_for_dest(dest_y, motions)} selects the
best right-leg motion; the target world coordinate is set to $(-0.5,\text{dest}_y,
0.11)$ (goal-plane center); the motion continues to play and the policy steers the
kick from the target term.

\paragraph{Target region $\to$ motion map.}
Built from \code{eval_motion_analysis.py} (5000 trials per motion$\times$target):

\begin{table}[h]
\centering\small
\caption{Target region to best motion (data-driven).}
\begin{tabular}{@{}lllll@{}}
\toprule
Target $y$ range & Best motion & Speed & Accuracy & Reason \\
\midrule
$[-1.5,-0.7)$ far left & \code{006_right} & 16.2\,m/s & \textbf{0.10\,m} & best far-corner accuracy \\
$[-0.7,-0.3)$ mid-left & \code{009_right} & 14.2\,m/s & 0.35\,m & best right-leg mid-left accuracy \\
$[-0.3,+0.3]$ center & \code{010_right} & \textbf{16.4\,m/s} & 0.11\,m & highest speed $+$ good center accuracy \\
$(+0.3,+0.7]$ mid-right & \code{005_right} & 12.1\,m/s & \textbf{0.006\,m} & best accuracy overall (sub-cm) \\
$(+0.7,+1.5]$ far right & \code{005_right} & 15.7\,m/s & 0.64\,m & most versatile motion \\
\bottomrule
\end{tabular}
\end{table}
All selected motions are right-leg kicks. Left-leg motions are unused because
they are slower (6--9\,m/s) and completely fail the far-left corner (0\% goal
rate).

\paragraph{Kick-frame map and zero-GK baseline.}
\code{analyze_kick_timing.py} confirms all motions contact the ball at frames
26--30 (0.52--0.60\,s) with zero std, which makes the \code{_should_lock()}
deadline reliable (e.g.\ \code{001_right} kicks at frame 27/260,
\code{005_right} at 26/230, \code{006_right} at 27/225). When \code{compete.py}
omits \code{--goalkeeper-api}, the GK is a \code{ZeroPolicy} (holds the
\code{HOME_KEYFRAME} crouch at $y=0$), giving a pure technical baseline.

\subsubsection{Motion performance analysis}\label{sec:sh-motion}

\textbf{Scale:} 10 motions $\times$ 5 targets $\times$ 5000 trials $=
\mathbf{250{,}000}$ trials, model Stage 6 (\code{model_20000.pt}), targets
$y=-1.50,-0.75,0.00,+0.75,+1.50$ (even split of the 3.0\,m goal).

\paragraph{Overall.} Mean goal rate $\mathbf{94\%}$ (9/10 motions at 100\%); kick
trigger rate $100\%$; right-leg mean ball speed $14.2$\,m/s ($\approx51$\,km/h);
left-leg $8.0$\,m/s ($\approx29$\,km/h); right-leg mean accuracy $0.42$\,m;
left-leg $0.26$\,m.

\begin{table}[h]
\centering\small
\caption{Right-leg motion performance (all 100\% goal rate).}
\begin{tabular}{@{}lcccccc@{}}
\toprule
Motion & Frames & Mean speed & Mean accuracy & Hit@0.3m & Best target & Best accuracy \\
\midrule
\code{001_right} & 260 & 13.59\,m/s & 0.340\,m & 40\% & $+0.75$ & 0.152\,m \\
\code{004_right} & 257 & 14.02\,m/s & 0.548\,m & 20\% & $+0.75$ & 0.197\,m \\
\code{005_right} & 230 & 13.96\,m/s & 0.402\,m & 20\% & \textbf{$+0.75$} & \textbf{0.006\,m} \\
\code{006_right} & 225 & 14.56\,m/s & 0.343\,m & 40\% & \textbf{$-1.50$} & \textbf{0.095\,m} \\
\code{009_right} & 205 & 13.73\,m/s & 0.454\,m & 40\% & $+0.00$ & 0.144\,m \\
\code{010_right} & 230 & 15.40\,m/s & 0.512\,m & 20\% & $+0.00$ & 0.112\,m \\
\bottomrule
\end{tabular}
\end{table}

\begin{table}[h]
\centering\small
\caption{Left-leg motion performance (80\% goal rate; 0\% only at $y=-1.50$).}
\begin{tabular}{@{}lcccc@{}}
\toprule
Motion & Frames & Mean speed & Mean accuracy & Hit@0.3m \\
\midrule
\code{002_left} & 214 & 7.27\,m/s & 0.245\,m & 80\% \\
\code{003_left} & 289 & 7.88\,m/s & 0.277\,m & 65\% \\
\code{007_left} & 229 & 8.42\,m/s & 0.266\,m & 60\% \\
\code{008_left} & 225 & 8.36\,m/s & 0.261\,m & 60\% \\
\bottomrule
\end{tabular}
\end{table}

\begin{table}[h]
\centering\small
\caption{Target $\times$ ball-speed matrix (m/s), right-leg motions.}
\begin{tabular}{@{}lccccc@{}}
\toprule
Motion & $y{=}{-}1.50$ & $y{=}{-}0.75$ & $y{=}{+}0.00$ & $y{=}{+}0.75$ & $y{=}{+}1.50$ \\
\midrule
\code{001_right} & 14.58 & 14.04 & 14.14 & 12.41 & 12.80 \\
\code{004_right} & 13.02 & 14.52 & 14.62 & 12.53 & 15.41 \\
\code{005_right} & 14.59 & 13.37 & 14.12 & 12.06 & \textbf{15.69} \\
\code{006_right} & \textbf{16.18} & 14.18 & 14.04 & 12.99 & 15.42 \\
\code{009_right} & 14.40 & 14.24 & 12.14 & 12.54 & 15.33 \\
\code{010_right} & 14.64 & 14.57 & \textbf{16.42} & 16.28 & 15.10 \\
\bottomrule
\end{tabular}
\end{table}

\paragraph{Key findings.} (1) \emph{Strong right/left asymmetry:} right-leg ball
speed is $\approx2\times$ the left, and the left leg cannot reach the far-left
corner -- likely because right-leg motions are more numerous/better in the data
or the reward structure favors the right leg. (2) \emph{Far-corner accuracy
degrades:} all motions lose accuracy at $y=\pm1.50$ (0.4--1.2\,m vs.\ 0.1--0.4\,m
at center), a reasonable mechanical limit. (3) \emph{Competition recommendation:}
use right-leg motions only, with the \code{gk-aware} strategy that reads the GK
during execution and shoots to the far side.

\subsubsection{Engineering decisions and lessons}\label{sec:sh-engineering}

\paragraph{Unified architecture.} All six stages share an identical LSTM
(128-64-32 $+$ 2$\times$128), differing only by environment config. This lets
checkpoints transfer seamlessly across stages (\code{--resume} loads
actor$+$critic$+$normalizer), concentrates effort on reward/environment design
rather than re-tuning networks, and speeds iteration (changing a reward weight
needs no network redesign).

\paragraph{Coordinate alignment is a fatal detail.} Stage 5 switching the
training frame fully to \code{compete.py}'s world coordinates
(robot $(4,0,0.8)$, yaw $\pi$) looks trivial but is decisive: it removes the
sim-to-sim domain gap. Checkpoints trained earlier in the training frame degrade
markedly in the compete scene.

\paragraph{LSTM state continuity.} Not calling \code{policy.reset()} on target
lock in \code{gk-aware} was a key lesson. An initial implementation reset on
\code{should_lock}, feeding the LSTM an empty hidden state at motion frame 16 -- a
combination never seen in training -- causing an OOD collapse. The fix updates
only the \code{dest_body} term and preserves LSTM continuity.

\paragraph{Data-driven motion selection.} The \code{_select_motion_for_dest()}
map is not heuristic but grounded in the 250K-trial measurements; each target
region's recommendation is backed by concrete speed/accuracy numbers, avoiding
subjective guessing.

\paragraph{Rule-based vs.\ adversarial training.} Phase 2 pits the shooter against
an unknown opponent GK. Rather than \emph{adversarial training} (which is
unstable when two policies learn online and would compound an already complex
six-stage curriculum), we chose a \emph{strong single agent $+$ rule-based target
selection}: first train a near-perfect GK-free shooter (99.96\% goal rate), then
add a \code{gk-aware} rule that reads the GK and shoots to the far side. ``Shoot
to the far side'' is a physically deterministic decision that needs no learning,
and given a shooter that scores 12--16\,m/s anywhere in the goal, this beats a
random target by a wide margin against a reactive GK. The trade-off is reduced
adaptivity to complex GK behavior (occlusion, tackles), but interpretability and
controllability are preferable against an unknown opponent to a possibly
under-trained adversarial policy.


\section{Goalkeeper}\label{part:keeper}

\subsection{Phase 1 Results}\label{sec:gk-results}

\paragraph{Evaluation setup.} We use the official Phase-1 goalkeeping
evaluation over three fixed seeds, with $50$ trials per seed ($150$ total). A
trial is counted as successful only if the ball does not enter the goal frame
before timeout.

\begin{table}[h]
\centering\small
\caption{Goalkeeper Phase-1 quantitative results.}
\label{tab:gk-phase1-summary}
\begin{tabular}{@{}lc@{}}
\toprule
Metric & Value \\
\midrule
Seed 42 block rate & $100.00\%$ \\
Seed 2810 block rate & $90.00\%$ \\
Seed 202606 block rate & $94.00\%$ \\
Mean block rate & $\mathbf{94.67\%}$ \\
Phase-1 score & $\mathbf{22.0/30}$ \\
\bottomrule
\end{tabular}
\end{table}

\paragraph{Comparison.} The evaluated goalkeeper remains clearly stronger than
the single repaired network ($81.7\%$) and the reference teacher from
Ren et al.~\cite{ren} ($\approx 72.5\%$ on our benchmark).

\paragraph{Analysis.} The result is strong and stable across seeds, though
high-corner saves remain the hardest cases. The spread from $90\%$ to $100\%$
across seeds suggests that most remaining misses are difficult boundary cases
rather than systematic failures of the routing design.

\subsection{Phase 2 Design \& Approach}

\paragraph{Phase-2 approach.} For the tournament phase, we build on the strong
Phase-1 goalkeeper and improve it through better repair, specialization,
routing, and deployment logic, so that it transfers more reliably to the
competition setting. The emphasis is on robustness and consistency under
tournament-style play.

\subsubsection{Task and the core difficulty}\label{sec:gk-task}

The goalkeeper must intercept a ball launched on a parabolic trajectory toward a
full $3\,\text{m}\times1.8\,\text{m}$ goal, controlling the G1 at $50$\,Hz. A
save is credited only if the ball never crosses the goal plane within the frame.
Crucially, the diving saves required to defend the corners necessarily terminate
with the robot on the ground.

This single property reshapes the entire problem. Any learning signal that
penalizes falling simultaneously penalizes the dives that produce saves, so the
conventional ``remain upright'' inductive bias is actively counterproductive; the
evaluation reflects this by not treating a fall as a failure. The consequence is
that \textbf{the required interception behavior is not reachable by reward-shaped
exploration} -- under a hand-designed reward, a humanoid does not stumble into a
committed, directional dive through random search. We verified this empirically:
trained from scratch with PPO~\cite{ppo} under a range of shaped rewards, the
policy reliably learns to stand and to satisfy the proxy terms, but its block
rate saturates between $3\%$ and $33\%$, with the high-ball corners frozen near
zero. A weak reward yields passive standing ($\approx 33\%$, barely above the
do-nothing baseline); a strong one induces an uncontrolled lunge that topples
without intercepting.

The central design question is therefore not ``which reward yields a goalkeeper''
but rather \textbf{how to inject a diving prior, and how to subsequently improve
upon it against the true objective}. The remainder of the method follows from
this premise.

A second observation establishes the target. A per-region diagnosis of the
Humanoid Goalkeeper of Ren et al.~\cite{ren} -- whose checkpoint serves as our
teacher -- shows that it blocks only $\approx 72.5\%$ on our benchmark, with
failures concentrated almost entirely in the high balls (low balls
$\approx 95\%$, the upper corners $\approx 40$--$60\%$). This is doubly
informative: competent behavior still leaves the corners largely undefended, and
it implies that \textbf{pure imitation of Ren et al.~\cite{ren} cannot exceed
$\approx 72.5\%$}. To do better, we must \emph{synthesize} high-ball saving
behavior that the teacher does not itself possess.

\subsubsection{Method: a staged pipeline culminating in a region mixture-of-experts}\label{sec:gk-method}

\paragraph{Prepare and active phases.} Our goalkeeper uses an explicit
\textbf{two-phase control design}. The first is the \emph{prepare} phase before
the shot becomes fully threatening. Its purpose is not to make an early save,
but to put the robot into a stable, ready-to-dive configuration: balanced over
the support polygon, crouched to a useful height, and aligned laterally so that
either-side launch remains available. This stage solves a practical failure
mode we repeatedly observed in single-phase designs: if the policy begins
reacting aggressively from the first moment, it often spends its control
authority on unnecessary sway, loses balance before the ball arrives, or enters
the interception too upright to generate a decisive dive.

The second is the \emph{active} phase after the ball trajectory is clearly
committed toward the goal. Here the objective changes completely: the
goalkeeper should stop preserving a neutral stance and instead execute the
region-appropriate interception as quickly and decisively as possible. In other
words, the design separates \emph{stability preparation} from \emph{save
execution}. This separation proved important because the two phases favor
different behaviors: before launch we want calm posture regulation, while after
launch we want explosive lateral and vertical motion even if the save ends on
the ground. Structuring the policy this way makes the keeper noticeably more
stable at shot onset and gives the later diving controller a much better state
from which to attack the ball.

Because the diving prior cannot be acquired from scratch and imitation of
Ren et al.~\cite{ren} is upper-bounded at $\approx 72.5\%$, no single training
procedure is sufficient. We adopt a \textbf{staged pipeline} in which each stage
overcomes the failure mode of the preceding one: distil a diving base
(Stage~A); \emph{prove} with an offline optimizer that the conceded balls are
physically saveable and fold those saves into one network (Stages~B--C);
\emph{specialize} a bounded corrective head per ball region (Stage~D); and route
the specialists with an analytic gate into the deployed mixture (Stage~E).

\paragraph{Stage A --- Distillation: injecting the diving prior.}

We behavior-clone the teacher policy of Ren et al.~\cite{ren} into a trainable
network of our own design, since the released teacher is inference-only under
the current RL framework~\cite{rudin} and cannot itself be optimized. The
procedure follows DAgger~\cite{dagger}: the student is rolled out, the teacher
labels the visited states, and the aggregated dataset is cloned. This yields a
competent, reactive \emph{diving} base policy with a $960$-D history-based
input. It does not surpass the teacher -- imitation cannot -- scoring
$\approx 74\%$ over $100$ trials and $\approx \mathbf{66.9\%}$
on the harder vectorized benchmark used in the rest of this part. It nevertheless
provides a trainable policy that already executes dives, the precondition for
every subsequent stage.

\paragraph{Stage B --- Repair oracle: proving the hard balls are saveable.}

To exceed the base policy we must save high balls that it concedes. We construct
a \textbf{per-scenario trajectory optimizer}: for the frozen base policy and a
fixed incoming ball, it searches -- using the cross-entropy method,
CEM~\cite{cem,icem}, a sampling-based optimizer -- an open-loop \emph{residual}
action sequence added to the base policy's actions. The search minimizes a dense
cost combining the true concede outcome with the distance from the nearest
blocking limb to the ball's predicted plane-crossing point. This optimizer
\textbf{demonstrates that the hard balls are physically saveable} at a
per-scenario save rate far above the base policy, establishing that control
\emph{authority} rather than kinematic reach was the limiting factor. The
optimizer is single-agent simulation optimization, not a learned policy -- it is
used as a \emph{data generator}, not as the deliverable.

\paragraph{Stage C --- Distilling the repaired saves into a single network.}

The repaired residuals are per-scenario and open-loop; they cannot be deployed
directly. We therefore \textbf{collect the repaired observation-action pairs
from trajectories that the oracle turns into successful saves, and distil them
back into a native MLP} with the same architecture as the base. The resulting
single repaired network reaches \textbf{81.7\%} -- a
$+14.8$-point gain over the $66.9\%$ base. This single reactive network is strong
but still pays the single-policy trade-off (improving one corner regresses
another); it becomes the \emph{shared base} that the region specialists of
Stage~D each refine. Two collection choices proved essential, each of which
silently degraded the result when violated:

\begin{itemize}
\item \emph{Demonstrations must be evaluation-matched.} Forcing a scenario after
the environment reset leaves stale ball-history frames in the observation,
training the policy on a distribution it never encounters at evaluation. The
scenario must be imposed \emph{through} the reset so that the observation history
is consistent with the deployed keeper observation/history layout.
\item \emph{Collect a stable save window, not the post-save fall.} A naive
collection learns the high-ball save together with the subsequent
topple-and-flail, so the deployed policy blocks but then loses balance after
contact. The oracle residual is faded back to the base policy after the save
window, and distillation keeps \textbf{only the frames around the keeper-plane
crossing instant} (a short pre/post window), so the student learns the
interception rather than the fall.
\end{itemize}

A \textbf{region-weighted, multi-GPU collection} path supports the same pipeline:
one repair worker per GPU, with region weights that oversample low/mid balls
(regions $4,5$ low, $0,1$ mid, $2,3$ upper) to keep the distilled policy from
overfitting to the dramatic high-ball jumps at the expense of routine balls.

\paragraph{Stage D --- Per-region ballistic-residual specialists.}

A single reactive network cannot simultaneously hold six high-fidelity region
dives -- this is a capacity limit, not a tuning one. We therefore \textbf{freeze
the repaired base and train one bounded residual head per ball region} (right and
left $\times$ low, mid, up). Each head is
zero-initialized -- so training starts exactly at the strong base -- and receives
explicit \emph{ballistic timing features}: the time and $(y,z)$ of the ball at
the keeper plane $x=0$ and at the goal plane $x=-0.5$, the incoming $v_x$, and the
ball speed. Trained with the stable save-window collection of Stage~C and a region
focus, each specialist masters its own balls
without the cross-region trade-off that caps the single network. Because the
residual is bounded and zero-initialized, a specialist can only \emph{refine} the
base dive for its region, never destabilize it.

\paragraph{Stage E --- Mixture of experts with an analytic ballistic gate.}

At inference the six specialists are combined into a single deployable bundle
and routed by an \textbf{analytic ballistic-crossing gate} -- no learned parameters and no
privileged information. From the raw ball position and velocity the gate
extrapolates, in closed form, the ball's state when it reaches the keeper plane:
the crossing time $t$, the lateral crossing $c_y$, the crossing height $c_z$, and
the descent rate $v_z$ at crossing. Thresholds map $c_z$ to a height band
($c_z<0.85$\,m $\to$ low, $c_z>1.35$\,m $\to$ up, else mid) and $c_y<0$ to the
left side, selecting one of the six experts. The region is \textbf{latched early}
-- as soon as the ball is unambiguously approaching ($v_x<-1$ and within
$5$\,m of the keeper) -- so the correct specialist drives the \emph{entire} dive
rather than a default expert mishandling the opening frames (which is fatal for
fast balls to the far side). Finally, an optional \textbf{left/right symmetry
mirror} serves a weak corner through its stronger mirrored expert, correcting
the teacher's residual left-bias. This mixture is
what lifts the keeper above both the single repaired network and the
$\approx 72.5\%$ teacher.

\paragraph{How the two-phase design fits the final system.} The deployed keeper
combines these ideas hierarchically. The \emph{prepare} phase maintains a
compact, balanced ready state while the ball is still in its pre-commitment
window, ensuring that the robot enters the save from a controllable posture
rather than from incidental motion. Once the incoming shot is sufficiently
clear, the controller switches into the \emph{active} phase and the analytic
ballistic gate selects the region specialist that should own the dive. This
division of labor is one of the main reasons the mixture works in practice:
specialists are asked to solve the hard interception itself, not the separate
problem of recovering from a poor pre-shot stance.

\subsubsection{Additional design evidence and negative results}\label{sec:gk-secondary}

Three findings around the mixture are reported honestly: one clean negative that
shaped the specialist training, one variance caution, and one in-progress path to
push past the current best.

\paragraph{On-policy ballistic-residual PPO (negative).} Before training the
specialists from oracle data, we tried the obvious alternative: freeze the base
and learn each bounded, zero-initialized residual head directly with
PPO~\cite{ppo} on the block reward. The runs \emph{stayed} near the frozen base
and produced no clear gain -- the conservative residual is safe but does not, on
its own, synthesize the hard-ball saves. This negative is exactly why the
Stage~D specialists are trained by distilling the repair-oracle demonstrations
rather than by on-policy reward: the productive signal lives in the oracle, not
in on-policy exploration of the residual.

\paragraph{Gate-threshold sweeping (variance caution).} The analytic gate has a
few scalar thresholds (the height bounds and the latch distance), and short-run
threshold sweeps can produce tempting candidates.
We found these unreliable: a candidate that looked best on a few hundred trials
regressed under a large-trial re-evaluation. We therefore fix the gate thresholds
and trust only the $16{,}384$-trial estimate, not short gate sweeps.

\paragraph{Failure mining and failure-replay PPO (in progress).} To push past the
current $\approx 93.8\%$, we also explored a pipeline that mines the mixture's
residual misses, replays those specific failing balls, and fine-tunes the
responsible specialist on them. It is included as the next-stage path but is
\textbf{not} the selected deployed policy, so we do not claim its gains
here.

\subsubsection{API server implementation and observation computation}\label{sec:gk-api}

For the cross-evaluation tournament the policy is served behind the standardized
competition API. The tournament driver sends the \textbf{same raw MuJoCo state}
to every team -- root pose and velocity, joint positions and velocities,
previous action, and ball position and velocity -- and our server turns that raw
state into exactly the observation tensor the policy was trained on. The match
between deployment-time observation construction and training-time observation
construction is what the entire mixture depends on, so we specify it
concretely.

\paragraph{Per-frame observation (96-D), computed server-side from raw state.}
On each control step the server builds one $96$-D frame by concatenating six
terms, each transformed into the goalkeeper's own pelvis frame with the exact
scalings used in training (Table~\ref{tab:gk-obs}). Only the ball position and
the keeper's own proprioception are used; no privileged or opponent state
enters the observation. The key non-obvious detail is that the reference frame
and the per-term scalings must match training exactly -- a wrong scale or a
world-frame ball vector can silently destroy the policy at deployment even
though the network itself loads without error.

\begin{table}[h]
\centering\small
\caption{Goalkeeper per-frame observation (96-D) built from raw state.}
\label{tab:gk-obs}
\begin{tabular}{@{}llp{8.3cm}@{}}
\toprule
Term & Dim & Computation from raw state \\
\midrule
Ball position in pelvis frame    & 3  & Relative ball position, rotated from world coordinates into the goalkeeper's local frame \\
Base angular velocity      & 3  & Root angular velocity rotated into the body frame and scaled consistently with training \\
Projected gravity & 3  & Gravity direction expressed in the body frame, giving an implicit orientation cue \\
Joint position offset     & 29 & Current joint angles relative to the goalkeeper default pose \\
Scaled joint velocity  & 29 & Joint velocities after the same normalization used during training \\
Previous action       & 29 & The action applied at the previous control step \\
\bottomrule
\end{tabular}
\end{table}

\paragraph{Term-major history stacking ($\to$ 960-D).} The actor consumes a
$10$-frame history, but \textbf{not} as a naive frame-major concatenation.
\emph{Term-major} stacking is used: all $10$ frames of ball position are placed
first, then all $10$ frames of base angular velocity, and so on per term,
yielding the $96\times 10 = 960$-D
vector the MLP expects. The server keeps a $10$-slot history buffer; on the first
frame after a reset it pre-fills all $10$ slots with the initial frame so the
very first action is well-defined.

\paragraph{Inference flow and mixture routing.} During play, the server computes
the $96$-D frame, updates the history buffer, stacks it to $960$ dimensions, and
runs policy inference. Because the analytic gate needs the raw ball state, the
server also tracks the ball position and velocity directly, extrapolates the
ballistic crossing point, latches the region early in the episode, and lets the
corresponding specialist control the rest of the dive. Reset clears the
per-episode latch and history state so each new trial starts cleanly.

\paragraph{Deployment status.} The main remaining integration requirement is to
ensure that the tournament deployment path preserves this same history-based
observation construction and mixture routing logic. A simplified single-frame
server would not serve the goalkeeper correctly, so the deployed system must
retain the same observation layout and routing behavior described above.

% ===========================================================================
\begin{thebibliography}{9}
% ===========================================================================
\bibitem{kong} J.~Kong, X.~Liu, Y.~Lin, J.~Han, S.~Schwertfeger, C.~Bai, and
X.~Li, ``Learning soccer skills for humanoid robots: A progressive
perception-action framework,'' 2026. arXiv:2602.05310.

\bibitem{ren} J.~Ren, J.~Long, T.~Huang, H.~Wang, Z.~Wang, F.~Jia, W.~Zhang,
J.~Wang, P.~Luo, and J.~Pang, ``Humanoid goalkeeper: Learning from
position-conditioned task-motion constraints,'' arXiv:2510.18002, 2025.

\bibitem{ppo} J.~Schulman, F.~Wolski, P.~Dhariwal, A.~Radford, and O.~Klimov,
``Proximal policy optimization algorithms,'' arXiv:1707.06347, 2017.

\bibitem{dagger} S.~Ross, G.~J.~Gordon, and J.~A.~Bagnell, ``A reduction of
imitation learning and structured prediction to no-regret online learning,'' in
\emph{Proc. AISTATS}, 2011, pp.~627--635.

\bibitem{icem} C.~Pinneri, S.~Sawant, S.~Blaes, J.~Achterhold, J.~Stueckler,
M.~Rolinek, and G.~Martius, ``Sample-efficient cross-entropy method for real-time
planning,'' in \emph{Proc. CoRL}, PMLR vol.~155, 2021, pp.~1049--1065.

\bibitem{cem} R.~Y.~Rubinstein and D.~P.~Kroese, \emph{The Cross-Entropy Method}.
New York: Springer, 2004.

\bibitem{mujoco} E.~Todorov, T.~Erez, and Y.~Tassa, ``MuJoCo: A physics engine
for model-based control,'' in \emph{Proc. IEEE/RSJ IROS}, 2012, pp.~5026--5033.

\bibitem{rudin} N.~Rudin, D.~Hoeller, P.~Reist, and M.~Hutter, ``Learning to walk
in minutes using massively parallel deep reinforcement learning,'' in
\emph{Proc. CoRL}, PMLR vol.~164, 2022, pp.~91--100.
\end{thebibliography}

\end{document}
